\documentclass[11pt]{article}
\usepackage{fontspec}
\usepackage[french,english]{babel}
\usepackage{amsmath}
\usepackage{amssymb}
\usepackage{amsthm}
\usepackage{mathrsfs}
\usepackage{tikz-cd}
\usetikzlibrary{babel}
\usepackage{enumitem}

\newcommand{\printedpage}[1]{\marginpar{\footnotesize p.~#1}}
\newcommand{\msp}[1]{\marginpar{\footnotesize [#1]}}
\newcommand{\D}{\mathbf{D}}
\newcommand{\cA}{\mathcal{A}}
\newcommand{\bDelta}{\boldsymbol{\Delta}}
\newcommand{\uHom}{\underline{\mathrm{Hom}}}

\begin{document}

\printedpage{}% no page number printed on this page
% head: (none)

\begin{center}
{\LARGE\textsf{THE DERIVATORS}}

\bigskip

\textbf{ALEXANDRE GROTHENDIECK}

\bigskip\bigskip

\textbf{Chapter IV}

\medskip

\textbf{Substantial diagrams}
\end{center}

\clearpage

\printedpage{}% no page number printed on this page
% head: (none)

\noindent This text was deciphered and transcribed in \LaTeXe\ by M. Künzer. It was edited by M. Künzer, J. Malgoire, and G. Maltsiniotis. The transcription is as faithful as possible to the manuscript. For the few obvious corrections, or rare comments by the editors, as well as for the original numbering of the manuscript pages, typewriter characters \texttt{[typewriter]} in brackets are used. A question mark in brackets means that one is not sure of the preceding word.

\noindent This edition is provisional. Remarks, comments and corrections are welcome. Send a message to:

\noindent\textbf{[e-mail address removed]}

\hfill G. Maltsiniotis\hspace*{2cm}

\clearpage

\printedpage{1}
% head: A. Grothendieck : Dérivateurs, Ch. IV, Édition des Universités de Montpellier II et Paris VII
% foot: Version 2 février 2002

\noindent\texttt{[There are no pages 1-72]}

\msp{page 73}

\subsection*{Construction of substantial diagrams for a derivator.}

\noindent Let $\D$ be a prederivator, $I$ in Diag, let us consider

\noindent $\varphi_I : \mathcal{D}(I) \longrightarrow \uHom(I^o,\cA)$,\ \ where $\cA=\D(e)$ is the fundamental category of coefficients. % [sic] \mathcal{D}(I) -> \mathbf{D}(I), elsewhere the prederivator is the bold D

The objects of $\uHom(I^o,\cA)$ are interpreted as \textit{diagrams of type $I^o$} in $\cA$. This language is especially convenient if $I$ is ordered. Let us denote, for $F\in\mathrm{Ob}\,\D(I)$, by $(F_i)_{i\in\mathrm{Ob}\,I}$ the associated $I^o$-diagram in $\cA$. By axiom Der 2, the functor
\[
\varphi : F \longmapsto (F_i)_{i\in\mathrm{Ob}\,I}
\]
is conservative. But it has no desire to be faithful, still less fully faithful. When one knows the diagram $(F_i)_i$, one is far from knowing $F$, not even up to nonunique isomorphism, I presume. It is by no means clear that if $F$, $F'$ in $\D(I)$ are such that the associated diagrams are isomorphic, they are themselves isomorphic.

To insist on this point, the objects of $\D(I)$ will be called the ``\textit{substantial diagrams} of type $I^o$ in $\cA$'' (or $\D$-substantial, if there is a risk of confusion). To support intuition and simplify the explanations, an object $F$ of $\D(I)$ ($\D$-coefficients on $I$) will be represented graphically by the corresponding ``weak'' or ``skeletal'' diagram -- one regards it
\msp{page 74}
as a sort of (pale) ``skeleton'' of the substantial diagram from which it comes. E.g. if $I=\bDelta_0 : 0\longrightarrow 1$ \texttt{[rather $\bDelta_1$]}, one denotes a substantial diagram of type $(\bDelta_0)^o$ \texttt{[rather $(\bDelta_1)^o$]} by the same graphical notation, as the ordinary diagram:
\[
F_0 \longleftarrow F_1.
\]
If
\[
I = \begin{tikzcd}
a \arrow[r] & c \\
0 \arrow[u] \arrow[r] & b, \arrow[u]
\end{tikzcd}
\]
the substantial diagrams of type $I$ are written
\[
\begin{tikzcd}
F_a \arrow[d] & F_c \arrow[l] \arrow[d] \\
F_0 & F_b. \arrow[l]
\end{tikzcd}
\]

\clearpage

\printedpage{2}
% head: A. Grothendieck : Dérivateurs, Ch. IV, Édition des Universités de Montpellier II et Paris VII
% foot: Version 2 février 2002

It will of course be necessary to specify each time whether this graphical notation designates a skeletal diagram, or a substantial one. Often one starts from a skeleton, and tries to ``flesh it out'' with a little flesh, to make it into a substantial diagram -- if possible, in a canonical way. This is a point of derivatic philosophy essential for practice. One will often proceed step by step, by ``substantifying'' subdiagrams corresponding to a strictly increasing sequence $I_0\subset I_1\subset\cdots\subset I_n=I$ of $I$, as best suits the particular situation.

\msp{page 75}

Here is a typical example, assuming $\D$ a \textit{pointed} derivator (the $\D(I)$ are pointed, and the $f^*$ take zero objects to zero objects). It is a matter of defining the ``cosuspension'' or ``loop space'' functor in $\cA$,
\[
\Omega : \cA \longrightarrow \cA.
\]
It is understood that for $M$ in $\cA$, $\Omega(M)$ must fit into a commutative square
\[
\begin{tikzcd}
e \arrow[d] & \Omega M \arrow[l] \arrow[d] \\
M & e, \arrow[l]
\end{tikzcd}
\tag*{(C)}
\]
where $e$ denotes the unit elements. This square however should be substantial, more precisely, be a substantial ``$\D$-\textit{cartesian}'' square. It being understood that an object of $\D(I)$ (substantial square in $\cA$) is called $\D$-cartesian, if it is in the essential image of $\D(I_1)$ under $i_{1*}$, where
\[
i_1 : I_1 \hookrightarrow I,
\]
$I_1$ being the open subcategory of $I$ defined by
\[
I_1 = \begin{tikzcd}
a \\
0 \arrow[u] \arrow[r] & b.
\end{tikzcd}
\]

\msp{page 76}

The restriction of the substantial square sought (C) to the subcategory $I_1$ is therefore a substantial diagram with skeleton
\[
\begin{tikzcd}
e \arrow[d] \\
M & e, \arrow[l]
\end{tikzcd}
\tag*{(Sq)}
\]
\printedpage{3}
% head: A. Grothendieck : Dérivateurs, Ch. IV, Édition des Universités de Montpellier II et Paris VII
% foot: Version 2 février 2002
from which it is obtained (up to canonical isomorphism) by applying $i_*$ to the latter, so as to complete it to $\begin{tikzcd}
e \arrow[d] & \Omega M \arrow[l, dotted] \arrow[d, dotted] \\
M & e \arrow[l]
\end{tikzcd}$. It therefore remains to describe a canonical substantiation of the skeleton (Sq), in terms of $M$. This leads us to look at the open subcategory
\[
I_0 = \{0\} :\quad
\begin{tikzcd}
{} \arrow[r, dotted] & {} \\
0 \arrow[u, dotted] \arrow[r, dotted] & {} \arrow[u, dotted]
\end{tikzcd}
\]
of $I_1$ (hence also of $I$) reduced to the initial object, and the inclusion
\[
I_0 \overset{i_0}{\hookrightarrow} I_1,
\]
and to substantify the skeleton (Sq) by $(i_0)_!(M)$.

One can say a posteriori that, if one knew the substantial square, there is, up to unique isomorphism, only one way of obtaining it as an $i_{1*}i_{0!}(?)$, with ? in $\D(I_0)\simeq\cA$, for $i_{0!}$ and $i_{1*}$ being fully faithful, the same holds for their composite.

\msp{page 77}

In fact, one sees that the functor
\[
i_{1*}i_{0!} :\ \cA\simeq\D(I_0) \longrightarrow \D(I)
\]
induces an equivalence of $\cA$ with the full subcategory of $\D(I)$ formed by the $\D$-cartesian squares whose fibers at $a$ and $b$ are zero. (The obvious quasi-inverse functor is $i^*$, where $i : U_0 \hookrightarrow I$ \texttt{[rather $I_0 \hookrightarrow I$]}.) This therefore means that in this case, the datum of the skeleton of a very particular nature ($F_a\simeq F_b\simeq 0_{\cA}$), plus the additional $\D$-cartesian constraint, on the square sought, determines the latter up to unique isomorphism in terms of its skeleton: there is a substantiation of the latter not only ``canonical'', but even unique up to unique isomorphism, because of the $\D$-cartesian constraint.

I would like to begin by looking at the still simpler case where $I=\bDelta_1=(0\longrightarrow 1)$, so the skeletons envisaged are simply the arrows
\[
F_1 \longrightarrow F_0 \qquad \text{in } \cA.
\]
The philosophy obliges us, for the most fundamental arrows in the theory of derivators $\D$ (arrows in $\cA$, or by extension in a $\D(I)\simeq\cA_{\D_I}$, fundamental category of the induced derivator $\D_I$) to ask ourselves whether there is not

\end{document}
